%%%%%%%%%%%%%%%%%%%%%%%%%%%%%%%%%%%%%%%%%%%%%%%%%%%%%%%%%%%%%%%%%%%%%%%%
\section{Proof of the upper bound}\label{sec:upper}
%%%%%%%%%%%%%%%%%%%%%%%%%%%%%%%%%%%%%%%%%%%%%%%%%%%%%%%%%%%%%%%%%%%%%%%%

\subsection{Gaussian multiplicative chaos}

Let $\Phi$ be a centered Gaussian field on a square whose covariance is $-\log|z-w|$ plus a bounded
function, such as the Gaussian free field. For $0<\theta<2$, the random measure
\[
e^{\theta \Phi(z)-\theta^2\E[\Phi(z)^2]/2}\,dz ,
\]
defined by smoothing $\Phi$ at scale $r$ and letting $r\to0$, is a nonzero random measure called
\emph{Gaussian multiplicative chaos}~\cite{RV}; subtracting half the variance makes the expected
density one. Its total mass over a fixed square has a heavy tail: the moment of order $p>0$ is
finite exactly when $p<4/\theta^2$~\cite[Theorem~2.11]{RV}. Intuitively, the mass can be carried by
a small region where $\Phi$ is unusually high. For $\Phi=h$ and $\theta=\gamma$ this measure is the
\emph{Liouville quantum gravity area}. \emph{Liouville quantum gravity} (LQG) with parameter $\gamma$
is the name for a random surface whose area element is, formally, $e^{\gamma h(z)}dz$; the chaos
construction makes this precise. It is used in Sections~\ref{sec:area} and~\ref{sec:limit}.

Two further facts are used. \emph{Kahane's convexity inequality}: if two centered Gaussian fields
$\Phi$ and $\Psi$ satisfy $\operatorname{Cov}(\Phi(z),\Phi(w))\le\operatorname{Cov}(\Psi(z),\Psi(w))$
for all $z,w$, then $\E[F(m_\Phi)]\le\E[F(m_\Psi)]$ for every convex
function $F$ of at most polynomial growth, where $m_\Phi$ and $m_\Psi$ are the total masses of the two
chaos measures with the same $\theta$~\cite[Theorem~2.1]{RV}. Intuitively, more correlation makes
the mass more variable. And for a field built from independent contributions of successive scales,
the masses at finer cutoffs form a martingale, uniformly integrable when
$\theta<2$~\cite[Theorem~2.5]{RV}.

\subsection{Uniform moments of the square counts}

For $\theta>0$ and a depth $j$, define the normalized sum over the $4^j$ squares of depth $j$,
\[
\mathcal M_{\theta,j}=4^{-j}\sum_{|S|=2^{-j}}\exp\Bigl(\theta h_S-\frac{\theta^2}2\operatorname{Var}(h_S)\Bigr).
\]
It is the total mass of the chaos measure of the piecewise constant field $\Phi_j(z)=h_{S_j(z)}$,
where $S_j(z)$ is the depth-$j$ square containing $z$.

\begin{lemma}\label{lem:moments}
If $0<\theta<2$ and $1\le p<4/\theta^2$, then $\sup_j\E[\mathcal M_{\theta,j}^p]<\infty$.
\end{lemma}

\begin{proof}[Proof sketch]
The covariance of the square averages gives, uniformly in $j$, $z$ and $w$,
\[
\operatorname{Cov}\bigl(\Phi_j(z),\Phi_j(w)\bigr)\le\log\frac1{|z-w|\vee2^{-j}}+C_0 .
\]
For comparison, take the field $\Psi_u$ with covariance $\int_0^ue^{-e^{2v}|z-w|^2/2}\,dv$, which has
independent scale increments; at $u=j\log2$ this covariance is $\log(1/(|z-w|\vee2^{-j}))+O(1)$.
Adding one independent centered normal variable of large fixed variance makes its covariance dominate that of $\Phi_j$ everywhere. Kahane's inequality with
$F(u)=u^p$ bounds $\E[\mathcal M_{\theta,j}^p]$ by the $p$-th moment of the comparison mass, times a
fixed lognormal moment. The comparison masses form a uniformly integrable martingale in $u$, so
each is the conditional expectation of the limiting chaos mass, and by Jensen's inequality its $p$-th
moment is at most that of the limit, which is finite because $p<4/\theta^2$. The passage from
Kahane's inequality for finite Gaussian vectors to these integrals is routine.
\end{proof}

\subsection{The upper bound for \texorpdfstring{$Q>2$}{Q>2}}

\begin{theorem}\label{thm:upper}
Let $Q>2$ and $x=4/\gamma^2$. Then $\E[N^p]<\infty$ for every real $1\le p<x$. Consequently
\[
\limsup_{n\to\infty}\frac{\log\Prob(N\ge n)}{\log n}\le-\frac4{\gamma^2}.
\]
\end{theorem}

\begin{proof}
Let $\mathcal N_j$ be the number of squares of depth $j$ that are active, whether or not they are reached
by the subdivision. Every square that is cut in the subdivision is active, and each cut adds
three squares, so
\[
N\le1+3\sum_{j\ge0}\mathcal N_j .
\]
For any $\theta>0$, the indicator of $\{e^{h_S/Q}2^{-j}>\varepsilon\}$ is at most
$(e^{h_S/Q}2^{-j}/\varepsilon)^{\theta Q}$, so
\[
\mathcal N_j\le\varepsilon^{-\theta Q}\,2^{-\theta Qj}\sum_{|S|=2^{-j}}e^{\theta h_S}.
\]
Since $\operatorname{Var}(h_S)=j\log2+O(1)$ uniformly, $e^{\theta h_S}\le
C\,2^{j\theta^2/2}\exp(\theta h_S-\theta^2\operatorname{Var}(h_S)/2)$, and therefore
\begin{equation}\label{eq:Ij}
\mathcal N_j\;\le\;C\,2^{\,j(2-\theta Q+\theta^2/2)}\,\mathcal M_{\theta,j}=C\,2^{\,j\psi(\theta)}\mathcal M_{\theta,j},
\end{equation}
with $\psi$ the quadratic~\eqref{eq:quadratic} and $C$ depending on $\theta$, $Q$ and $\varepsilon$.

The roots of $\psi$ are $\gamma$ and $4/\gamma$, and $\psi<0$ strictly between them. Fix $p<x$. Since
$\gamma<2$ we can choose $\theta$ slightly above $\gamma$ with $\theta<2$ and $p<4/\theta^2$; then
$\psi(\theta)<0$. By Minkowski's inequality, \eqref{eq:Ij} and Lemma~\ref{lem:moments},
\[
\|N\|_p\le1+3\sum_{j\ge0}\|\mathcal N_j\|_p\le1+C\sum_{j\ge0}2^{j\psi(\theta)}\sup_i\|\mathcal M_{\theta,i}\|_p<\infty .
\]
Markov's inequality gives $\Prob(N\ge n)\le\E[N^p]\,n^{-p}$. Letting $p$ increase to $x$ proves the
bound on the limit superior.
\end{proof}

\subsection{The upper bound at the critical endpoint}

At $Q=2$ we have $x=1$, so Minkowski's inequality is not available for the needed moments below
one. Instead:

\begin{proposition}\label{prop:critupper}
At $Q=2$ and every fixed $\varepsilon>0$, $\E[N]<\infty$. Hence $\Prob(N\ge n)\le\E[N]/n$.
\end{proposition}

\begin{proof}[Proof sketch]
Let $\mathcal R_j$ count the squares of depth $j$ that are active together with all their ancestors.
Then $N\le1+3\sum_j\mathcal R_j$, and $\E[\mathcal R_j]\le4^j\sup\Prob(\mathcal A_j)$, where $\mathcal A_j$ is the event that all squares of
a fixed chain $S_0\supset\cdots\supset S_j$ are active, that is $h_{S_l}>2l\log2+2\log\varepsilon$.

Change the law by the density $\exp(2h_{S_j}-2\operatorname{Var}(h_{S_j}))$, an exponential tilt.
This gives
\[
4^j\Prob(\mathcal A_j)\le C\,\E^*\bigl[e^{-2(h_{S_j}-2j\log2)}\,;\mathcal A_j\bigr],
\]
where $\E^*$ is the tilted expectation. Under the tilt, the field's averages over circles
centered in $S_j$ with radii $r_02^{-l}$ form a Brownian motion in the variable $l\log2$, with
drift exactly $2$, plus an independent normal starting value; the square averages differ from the
circle averages by Gaussian amounts of bounded variance. The threshold also has slope $2$. So,
relative to the threshold, the averages along the chain of ancestors form a Brownian motion
$\mathcal W$ without drift, which must stay above a barrier for a time $j\log2$. Up to the bounded
errors, $e^{-2(h_{S_j}-2j\log2)}$ is $e^{-2\mathcal W_{j\log2}}$.
Allow the errors to reach $O(\sqrt{\log j})$, at the cost of an event of probability
$O(j^{-M})$ for any fixed $M$. The reflection principle gives
\[
\E\bigl[e^{-2\mathcal W_{j\log2}};\min_{u\le j\log2}\mathcal W_u\ge-b'\bigr]\le C\,b'\,e^{2b'}\,(j\log2)^{-3/2}\qquad(b'\ge1).
\]
With $b'$ of order $\sqrt{\log j}$, this yields $\E[\mathcal R_j]\le j^{-3/2+o(1)}$, which is summable.
\end{proof}

Intuitively, at the critical point the factor $4$ in the number of squares per level exactly
balances the cost of staying active, and what remains is the probability $j^{-3/2}$ that a
driftless walk stays above a barrier and ends near it.
